% !TeX root = ../组合数学笔记.tex
% 由主文档通过 \input 载入；请编译主文档。
\setcounter{chapter}{1} % 保留原书章号。
\chapter{代表性习题解答}
\label{app:exercises}

本附录选择各章两道习题，给出关键论证或完整计算。未收录的练习应先按其计数对象、边界条件和目标结论独立尝试；计算答案相同并不意味着两种建模必然等价。

\section{基本方法与精确计数}

\subsection{第\ref{ch:objects}章}

习题\ref{ex:objects-1}：先选三名代表再选召集人，有 $\binom53\cdot3=30$ 种；先选召集人，再从其余四人中选两名普通代表，有 $5\binom42=30$ 种。两种过程都唯一指定同一个“委员会加负责人”对象。

习题\ref{ex:objects-2}：纤维两两不交且覆盖 $X$，故 $|X|=\sum_{y\in Y}|f^{-1}(\{y\})|\le d|Y|$。等号当且仅当每个纤维均恰有 $d$ 个元素；若存在一个严格小于 $d$ 的纤维，其余项不超过 $d$，总和便严格小于。

\subsection{第\ref{ch:elementary-counting}章}

习题\ref{ex:elementary-counting-1}：$\mathrm{MISSISSIPPI}$ 中 $I,S$ 各四个，$P$ 两个，$M$ 一个，故不同排列数为 $11!/(4!4!2!)=34650$。每个无副本标签的排列恰对应 $4!4!2!$ 个临时标号排列。

习题\ref{ex:elementary-counting-3}：令 $y_i=x_i-2$，得到四个非负整数之和为 $4$，故方案数为 $\binom73=35$。平移是双射，不改变各盒名称。

\subsection{第\ref{ch:identities}章}

习题\ref{ex:identities-2}：对 $n\ge2$，分解 $k^2=k(k-1)+k$，得到 $n(n-1)2^{n-2}+n2^{n-1}=n(n+1)2^{n-2}$。$n=0$ 时为零，$n=1$ 时为一；后一值也符合上述最终表达式。

习题\ref{ex:identities-3}：分成三段，路径数为 $\binom21\binom32\binom42=36$。若两个必经点互不可比，即一个横坐标更大但纵坐标更小，只向右或向上移动的路径不可能同时经过它们，计数为零。

\subsection{第\ref{ch:pigeonhole}章}

习题\ref{ex:pigeonhole-3}：取序列 $3,2,1,6,5,4,9,8,7$。递增子序列每个递减块至多取一项，所以长度至多三；递减子序列不能跨越数值整体增大的块，也至多长三。

习题\ref{ex:pigeonhole-4}：原操作保持所有数的总和，所以最后是 $1+\cdots+10=55$。改成 $a+b+1$ 后每步总和增加一，而数的个数从十减至一需要九步，最终为 $64$；两种结论都不依赖选取次序。

\subsection{第\ref{ch:inclusion-exclusion}章}

习题\ref{ex:inclusion-exclusion-1}：由错排递推或容斥得 $D_6=265$。恰有两个固定点时，先选固定位置，再让其余四个位置错排，得 $\binom62D_4=15\cdot9=135$。

习题\ref{ex:inclusion-exclusion-3}：容斥得到 $\binom{13}{3}-4\binom93+6\binom53=10$。也可令 $y_i=3-x_i$，得到 $\sum_i y_i=2$；各 $y_i\le3$ 自动满足，隔板法给出 $\binom53=10$。第二种方法利用了总数接近上界之和的特殊结构。

\subsection{第\ref{ch:recurrences}章}

习题\ref{ex:recurrences-3}：令 $a_0=1,a_1=2,a_2=4$。对 $n\ge3$，合法串按尾段分成以 $0$、$01$、$011$ 结尾的三类，删去该尾段分别得到长度 $n-1,n-2,n-3$ 的合法串。因此 $a_n=a_{n-1}+a_{n-2}+a_{n-3}$，继而 $a_3=7,a_4=13,a_5=24$。

习题\ref{ex:recurrences-4}：特征多项式为 $(\lambda-2)^2$，设 $a_n=(A+Bn)2^n$。初值得 $A=1,A+B=2$，故 $a_n=(n+1)2^n$。代入递推两侧均得 $(n+1)2^n$，完成核对。

\section{生成函数与经典结构}

\subsection{第\ref{ch:ogf}章}

习题\ref{ex:ogf-1}：生成函数为 $\frac1{1-x^2}\frac{x}{1-x}(1+x+x^2)$。固定 $x_3=0,1,2$，在 $n=5$ 时分别有 $3,2,2$ 个可选的非负偶数 $x_1$，且 $x_2$ 随即确定为正整数，故共七个解。

习题\ref{ex:ogf-3}：设三种面值的枚数为 $a,b,c$，则 $a+b+c=4,a+2b+5c=10$，相减得 $b+4c=6$。非负条件只允许 $c=1,b=2,a=1$，故只有一种。相应系数为 $[x^{10}y^4]((1-yx)(1-yx^2)(1-yx^5))^{-1}$。

\subsection{第\ref{ch:egf}章}

习题\ref{ex:egf-1}：递推得到 $S(5,2)=15,S(5,3)=25$；两个有名称非空盒子的分配数为 $2^5-2=30$，除去块的两种命名顺序，得到 $S(5,2)=15$。

习题\ref{ex:egf-3}：组件大小只许为正偶数，生成函数为 $\exp(\sum_{j\ge1}x^{2j}/(2j)!)$。四个标签或组成一个四元块，或组成两个二元块，方案数为 $1+\binom42/2=4$。

\subsection{第\ref{ch:classical-structures}章}

习题\ref{ex:classical-structures-2}：凸 $(n+2)$ 边形对应 $C_n$，故七边形对应 $n=5$，三角剖分数为 $C_5=\binom{10}{5}/6=42$。

习题\ref{ex:classical-structures-4}：将 $8$ 拆成八个 $1$，将 $6$ 拆成两个 $3$，再保留原来的 $3$ 和 $1$，得到三个 $3$ 与九个 $1$。总和仍为 $18$，重数的二进制合并又可恢复原分拆。

\section{图、对称与偏序}

\subsection{第\ref{ch:graphs-trees}章}

习题\ref{ex:graphs-trees-2}：依次选出的最小未出现在剩余序列中的标签为 $1,3,2$，连边为 $12,23,24$，最后连 $45$。因此边集为 $\{12,23,24,45\}$，度数依次为 $(1,3,1,2,1)$。

习题\ref{ex:graphs-trees-3}：Pr\"ufer 序列长度为三，其中标签 $1$ 出现两次、$2$ 出现一次，其他不出现，故有 $3!/(2!1!)=3$ 棵树。

\subsection{第\ref{ch:matching-coloring}章}

习题\ref{ex:matching-coloring-3}：四圈的染色多项式为 $(q-1)^4+(q-1)$，所以 $q=3$ 时有 $16+2=18$ 种。颜色已有名称，不再除以颜色排列数。

习题\ref{ex:matching-coloring-4}：顶点数至少三时，边数上界给出平均度数至多 $6-12/n<6$，故有度数至多五的顶点；一、二点图直接成立。删除该顶点后仍可平面，对剩余图按顶点数归纳作六色染色；加回时其至多五个邻点最多禁用五种颜色，尚有一种可选。这证明六色上界，不声称最少颜色数总为六。

\subsection{第\ref{ch:symmetry}章}

习题\ref{ex:symmetry-1}：五点二色旋转固定点总和为 $2^5+4\cdot2=40$，除以五得八类。五个反射各固定 $2^3=8$ 种图案，全部十个变换平均为 $(40+40)/10=8$，所以此例加入反射后数量不变。

习题\ref{ex:symmetry-3}：恒等旋转固定十五个染色，半周旋转固定三对相对红点，其他旋转没有合法固定染色，故得到三类。也可直接按两红点的较短圆周距离为一、二、三分类。

\subsection{第\ref{ch:posets}章}

习题\ref{ex:posets-1}：极小元为 $2,3$，没有最小元；极大元与最大元都为 $6$。最小元必须同时整除集合中的每个元素，不能把任一极小元都称为最小元。

习题\ref{ex:posets-4}：可分为链 $1<2<4<12$ 与 $3<6$（此处小于表示严格整除）。$\{4,6\}$ 为二元反链，任何链划分至少两条；上述构造达到下界，故最少为二。

\section{极值、概率、设计与算法}

\subsection{第\ref{ch:extremal-ramsey}章}

习题\ref{ex:extremal-ramsey-1}：上界为 $\lfloor49/4\rfloor=12$，完全二分图 $K_{3,4}$ 达到它。

习题\ref{ex:extremal-ramsey-3}：$R(3,4)\le R(2,4)+R(3,3)=4+6=10$。这只证明十点足够；要证明最小值恰为十，还需在九点上构造不含红三角形且不含蓝四点团的染色，而递推并未提供这样的构造。

\subsection{第\ref{ch:probabilistic}章}

习题\ref{ex:probabilistic-1}：每个位置成为固定点的概率为 $1/n$，所以 $\mathbb E X=\sum_{i=1}^n1/n=1$。第一矩法的充分条件要求严格小于一，此处不满足；尤其 $n=1$ 时确实没有错排。错排存在性需要别的论证或更细的信息。

习题\ref{ex:probabilistic-3}：函数 $f(p)=np-mp^2$ 是开口向下的二次函数，无约束最大点为 $p=n/(2m)$。若 $2m\ge n$，最大下界为 $n^2/(4m)$；若 $2m<n$，在区间 $[0,1]$ 上最大点为 $p=1$，下界为 $n-m$。实际独立集大小是整数，可将所得实数下界向上取整。

\subsection{第\ref{ch:designs-codes}章}

习题\ref{ex:designs-codes-1}：$r=\lambda(v-1)/(k-1)=4$，$b=vr/k=12$。这只检查了两项必要计数恒等式；若要证明存在，还须给出区组或引用满足全部条件的构造定理。

习题\ref{ex:designs-codes-4}：半径一球大小为 $1+7=8$，所以 $|C|\le2^7/8=16$。由七个不同非零三维二元列组成校验矩阵的零空间给出大小十六、最小距离三的线性码，恰好达到上界。

\subsection{第\ref{ch:algorithms}章}

习题\ref{ex:algorithms-1}：$2413$ 的各位较小未用标签数为 $1,2,0,0$，故排名为 $1\cdot3!+2\cdot2!=10$。编号十三依次给出商 $2,0,1,0$，在剩余有序标签中取对应位置，恢复 $3142$。

习题\ref{ex:algorithms-3}：若第一人取任务一，第二人只能取二，第三人取三；若第一人取二，第二人只能取三，第三人取一。因此完整分配为 $(1,2,3)$ 与 $(2,3,1)$ 两种。子集动态规划在最终状态 $\{1,2,3\}$ 处应得到二，小规模列举在此提供独立核对。

